\documentclass[11pt]{article}
\usepackage{fontspec}
\usepackage[french,english]{babel}
\usepackage{amsmath}
\usepackage{amssymb}
\usepackage{amsthm}
\usepackage{mathrsfs}
\usepackage{tikz-cd}
\usepackage{enumitem}

\newcommand{\printedpage}[1]{\marginpar{\footnotesize p.~#1}}
\newcommand{\msp}[1]{\marginpar{\footnotesize [#1]}}
\newcommand{\uHom}{\underline{\mathrm{Hom}}}
\newcommand{\Ind}{\mathrm{Ind}}
\newcommand{\Hom}{\mathrm{Hom}}
\newcommand{\card}{\mathrm{card}}
\newcommand{\Fl}{\mathrm{Fl}}
\newcommand{\Ob}{\mathrm{Ob}}
\newcommand{\Filt}{\mathrm{Filt}}
\newcommand{\hookxrightarrow}[2][]{\lhook\joinrel\xrightarrow[#1]{#2}}

\begin{document}

\setcounter{footnote}{11}

\printedpage{14}
% head: A. Grothendieck : Dérivateurs, Ch. XVIII, Édition des Universités de Montpellier II et Paris VII

whence
\[
\Fl\,S/x \;\simeq\; \varinjlim_{i} \Fl\,S/x_i \;,
\]
hence
\[
\card\Fl(S/x) \;\leqslant\; (\underbrace{\card I}_{\leqslant\pi})\underbrace{(\sup_{i}\card\Fl(S/x_i))}_{\leqslant\pi} \;\leqslant \pi^2 = \pi \;,
\]

\hfill q.e.d.

\msp{page 76}
% margin: hand-drawn bracket mark in the left margin beside the next two lines
Let us prove c) (\footnote{\textit{Botched proof}, only works if $\pi^{\pi_0} = \pi$\,!!}), i.e. that for every $\pi' > \pi \geqslant \pi_0$ and every $x \in \Ob\Filt^{\pi'}$, one can write
\[
x \;=\; \varinjlim_{I} x_i \qquad I \text{ large relative to } \pi\text{, the } x_i \text{ in } \Filt^{\pi}\text{, } \card I \leqslant \pi'^{\pi}. \tag*{(4.3.6.5)}
\]
% margin: hand-drawn bracket mark in the left margin beside the next three lines
The proof goes as in loc. cit. (p. 151-152), slightly changing the definition of $I$, as denoting the set of full subcategories $i$ of $S/x$ which are filtering, of cardinality $\leqslant \pi$ \textit{and moreover large relative to} $\pi_0$ (\footnote{It is not certain that $I$ is large relative to $\pi$, unless one assumes $\pi^{\pi_0} = \pi$\,!}). (We may suppose $S/x$ filtering by enlarging $S$ so that $S$ is stable in $\mathcal{M}$ under small finite $\varinjlim$, hence the categories $S/x$ too, which implies that they are filtering. It is \texttt{[incomplete sentence]})

We have thus proved PTA $\Longrightarrow$ c).

\textit{Let us prove c) $\Longrightarrow$ d).} Let
\[
\pi_1 \;=\; 2^c \qquad \text{with} \qquad c \;\geqslant\; \pi_0 \;.
\]
We then know that
\[
\Filt^{\pi_1}\mathcal{M} \;=\; \mathcal{M}_{\pi_1} \;\overset{\text{def}}{=}\; \text{full subcategory of $\mathcal{M}$ of the $\pi_1$-accessible objects of $\mathcal{M}$}
\]
(loc. cit. 9.16), which implies in particular that $\mathcal{M}_{\pi_1}$ is stable under $\varinjlim$ of cardinality $\leqslant \pi_1$, by virtue of loc. cit. 9.9 (page 145). Moreover,
\msp{page 77}
by loc. cit. 9.18 b) (p. 156), we have
\[
\Ind_{\pi_1}(\mathcal{M}_{\pi_1}) \xrightarrow{\;\sim\;} \mathcal{M} \;,
\]
which proves d).

\textit{Let us prove d) $\Longrightarrow$ PTA.} By loc. cit. 9.18 a) (p. 156), we see that the objects of $\mathcal{M}_{\pi_1}$ are $\pi_1$-accessible in $\Ind_{\pi_1}(\mathcal{M}_{\pi_1})$, on the other hand it is clear that $\mathcal{M}_{\pi_1}$ is a strictly generating full subcategory of $\Ind_{\pi_1}(\mathcal{M}_{\pi_1})$, so it remains to see that $\Ind_{\pi_1}(\mathcal{M}_{\pi_1})$ is stable under small $\varinjlim$ , i.e. the

\textbf{Lemma 4.3.7.} \textit{Let $\mathcal{M}_\pi$ be a $\mathfrak{U}$-category stable under direct limits of cardinality $\leqslant \pi$, where $\pi$ is a given infinite cardinal. Then} $\Ind_\pi\mathcal{M}_\pi$ \textit{is stable under small direct limits. Moreover, $\mathcal{M} = \Ind_\pi\mathcal{M}_\pi$, equipped with $\mathcal{M}_\pi \longrightarrow \mathcal{M}$, is a solution of a $2$-universal problem, namely
\printedpage{15}
% head: A. Grothendieck : Dérivateurs, Ch. XVIII, Édition des Universités de Montpellier II et Paris VII
that for every category $\mathcal{N}$ stable under small direct limits, the restriction functor induces an equivalence}
\[
\uHom_{!}(\mathcal{M},\mathcal{N}) \xrightarrow{\;\sim\;} \uHom_{!\,\pi}(\mathcal{M}_\pi,\mathcal{N}) \;, \tag*{(4.3.7.1)}
\]
\textit{where the sign ! (resp. $!\,\pi$) denotes the full subcategory of the category} $\uHom$\textit{, formed by the functors which commute with small} $\varinjlim$ \textit{(resp. with direct limits of cardinality $\leqslant \pi$).}

\msp{page 78}
Let us first prove that the $\varinjlim$ of cardinality $\leqslant \pi$ \textit{in} $\mathcal{M}_\pi$ are also $\varinjlim$ \textit{in} $\mathcal{M} = \Ind_\pi\mathcal{M}_\pi$ (\footnote{But \textit{beware}, $\mathcal{M}_\pi \longrightarrow \Ind(\mathcal{M}_\pi)$ does \textit{not} commute with these $\varinjlim$ !}). We thus have
\[
x \xleftarrow{\;\sim\;} \varinjlim_{I} x_i \qquad \textit{in } \mathcal{M}_\pi \;,
\]
i.e.
\[
\Hom(x,y) \xrightarrow{\;\sim\;} \varprojlim_{I} \Hom(x_i,y) \qquad \text{if } y \in \Ob\mathcal{M}_\pi \;, \tag*{$(*)$}
\]
let us prove that this remains true for $y$ in $\mathcal{M}$, i.e. $y$ of the form
\[
y \;=\; \varinjlim_{J} y_j \;, \qquad J \text{ large relative to } \pi\text{, the } y_j \text{ in } \mathcal{M}_\pi.
\]
We have indeed
\[
\begin{aligned}
\Hom(x,y) \quad &\simeq \quad \varinjlim_{J} \Hom(x,y_j)\\
&\overset{(*)}{\simeq} \quad \varinjlim_{J}\, \varprojlim_{I} \Hom(x_i,y_j)\\
&\overset{\text{SGA 4 I 9.8}}{\simeq} \quad \varprojlim_{I}\, \varinjlim_{J} \Hom(x_i,y_j)\\
&\simeq \quad \varprojlim_{I} \Hom(x_i,y) \;,
\end{aligned}
\]
altogether
\[
\Hom(x,y) \;\simeq\; \varprojlim_{I} \Hom(x_i,y),
\]
q.e.d.

To prove that the small $\varinjlim$ are representable in $\mathcal{M} = \Ind_\pi(\mathcal{M}_\pi)$, as they are in $\Ind(\mathcal{M}_\pi)$ (SGA 4 I 8.9.5 b, we use the hypothesis that the \textit{finite} $\varinjlim$ are representable in $\mathcal{M}_\pi$), it remains to see that $\Ind_\pi(\mathcal{M}_\pi)$ is stable in $\mathcal{M}_\pi$ under these limits.

\msp{page 79}

\section*{4.4\quad Study of $\Ind_\pi(C)$ (for $C$ a small category stable under $\varinjlim$ of cardinality $\leqslant \pi$)}

\textbf{4.4.1.} In the end, the proof of Lemma 4.3.7 is not so obvious as I anticipated, and I must devote some technical preliminaries to it, which I am going to collect here. Throughout what follows, $C$ represents a $\mathfrak{U}$-category equivalent to a small category.

\printedpage{16}
% head: A. Grothendieck : Dérivateurs, Ch. XVIII, Édition des Universités de Montpellier II et Paris VII

\textbf{Proposition 4.4.2.} \textit{Let $F \in C^{\wedge}$ be an ind-representable functor, i.e. such that $C/F$ is filtering. Let, on the other hand, $i \longmapsto x_i$ be a functor $I \longrightarrow C$, where $I$ is a small filtering category, and let}
\[
\alpha \;:\; x \;=\; \varinjlim_{I}{}_{C^{\wedge}} x_i \longrightarrow F \qquad \text{(\footnotemark)} \tag*{(4.4.2.1)}
\]
\footnotetext{Giving $\alpha$ is also equivalent to giving a homomorphism of ind-objects $(x_i)_{i\in I} \longrightarrow F$, $F$ being considered as a (generalized) ind-object indexed by $C/F$.}
\textit{an inverse system of arrows $x_i \longrightarrow F$ (}\textbf{N.B.} \textit{the direct limit $x$ is taken in $C^{\wedge}$, and is computed argument by argument), which is equivalent to a factorization $\varphi_\alpha : I \longrightarrow C/F$ of $I \longrightarrow C$ as}
\[
I \xrightarrow{\;\varphi_\alpha\;} C/F \hookxrightarrow[\psi]{\textit{can.}} C \;.
\]
\textit{That said, for $\alpha$ (4.4.2.1) to be an isomorphism in} $\Ind'(C) \subset C^{\wedge}$ \textit{(full subcategory of ind-representable functors on $C$), it is necessary and sufficient that $\varphi_\alpha$ be} cofinal\textit{, or what amounts to the same thing, $I$ being filtering, that it satisfy the two usual conditions:}
\begin{itemize}[leftmargin=2.5em]
\item[F\,1)] $\forall\, \xi \in \Ob C/F$, $\exists\, i \in \Ob I$ \textit{and an arrow} $\xi \longrightarrow \varphi_\alpha(i)$.
\item[F\,2)] $\forall\, \xi \in \Ob C/F$, \textit{and} $i \in \Ob I$, \textit{and every double arrow}

\msp{page 80}
\[
\xi \rightrightarrows \varphi_\alpha(i) \;,
\]
\textit{there exists an arrow $i \xrightarrow{\;u\;} j$ in $I$ such that $\varphi_\alpha(u)$ equalizes the double arrow.}
\end{itemize}
\textsc{Proof}. If $\varphi_\alpha$ is cofinal, we conclude from this by definition that $\varinjlim \varphi_\alpha \longrightarrow \varinjlim \psi\varphi_\alpha$ is an isomorphism, now this arrow is none other than $\alpha$. Conversely, suppose that this arrow is an isomorphism, hence $C/F \xleftarrow{\;\sim\;} C/\varinjlim x_i$, and it remains to see that the canonical functor
\[
I \;\longrightarrow\; C/\varinjlim{}_{C^{\wedge}} x_i
\]
is cofinal, i.e. satisfies F\,1) and F\,2). The verification is $\pm$ tautological, using the fact that for every $y$ in $C$, we have
\[
\Hom(y,\, \varinjlim{}_{C^{\wedge}} x_i) \;\xleftarrow{\;\sim\;}\; \varinjlim{}_{\mathrm{Ens}}\Hom(y,x_i),
\]
and the ``calculation'' of filtering direct limits in Ens.

% margin: vertical bar in the left margin beside Corollaire 4.4.3
\textbf{Corollary 4.4.3.} \textit{Let $\pi$ be an infinite ordinal. For $F$ to be in} $\Ind'_\pi(C)$ \textit{(essential image of} $\Ind_\pi(C)$ \textit{in $C^{\wedge}$), it is necessary and sufficient that there exist an ordered category $I$ large relative to $\pi$, and a cofinal functor $I \longrightarrow C/F$.}

We must therefore give conditions, on a category $J$ $(= C/F)$, for there to exist an ordered category $I$ large relative to $\pi$, and a cofinal functor $I \longrightarrow J$.

\msp{page 81}
Such a category will be called \textit{large relative to} $\pi$. (It is immediate that in the case where $J$ is itself ordered, this terminology is consistent with that already adopted for ordered sets, i.e. that if there exists an increasing cofinal map $\varphi : I \longrightarrow J$, with $I$ large relative to $\pi$, then $J$ is itself large relative to $\pi$.) Such a category is
\printedpage{17}
% head: A. Grothendieck : Dérivateurs, Ch. XVIII, Édition des Universités de Montpellier II et Paris VII
automatically filtering (SGA 4 I 8.1.3 b)). Let us also note that if we have a cofinal functor $J \longrightarrow J'$, if $J$ is large relative to $\pi$, $J'$ is too (by transitivity of cofinality).

\textbf{Proposition 4.4.4.} \textit{Let $J$ be a $\mathfrak{U}$-category, $\pi$ an infinite cardinal. Equivalent conditions:}
\begin{itemize}[leftmargin=4em]
\item[\textit{a)}] \textit{$J$ is large relative to $\pi$, i.e. there exists an ordered set large relative to $\pi$ and a cofinal functor $\varphi : I \longrightarrow J$.}
\item[\textit{b)}] \textit{$J$ satisfies the following two conditions:}
\begin{itemize}[leftmargin=4em]
\item[$\mathrm{PF}_{\pi}\,1$] $\mathrm{Ord}\,J$ \textit{is large relative to $\pi$, i.e. for every family $(y_\alpha)_{\alpha\in A}$ of objects of $J$, with} $\card A \leqslant \pi$\textit{, there exists an upper bound $y$ in $J$ (i.e. a $y$ and arrows $y_\alpha \xrightarrow{\;v_\alpha\;} y$).}
\item[$\mathrm{PF}_{\pi}\,2$] \textit{For $y_0$, $y_1$ in $J$ and every set $(u_\alpha)_{\alpha\in A}$ of arrows $y_0 \xrightarrow{\;u_\alpha\;} y_1$ with} $\card A \leqslant \pi$\textit{, there exists an arrow $y_1 \xrightarrow{\;u'\;} y_2$ which equalizes them, i.e. such that $u'u_\alpha = u'u_\beta \;\forall\, \alpha,\beta \in A$. (It suffices to require this when} $\card A \geqslant 2$.\textit{)}
\end{itemize}
\item[\textit{c)}] \textit{For every subset $A$ of} $\Fl J$\textit{, with} $\card A \leqslant \pi$\textit{, exists a subcategory $J'$ of $J$, with} $\card\Fl J' \leqslant \pi$\textit{, such that $A \subset \Fl J'$ and $J'$ has a final object.}
\item[\textit{d)}] \textit{Let $I$ be the set of subcategories $J'$ of $J$ such that} $\card J' \leqslant \pi$ \textit{and having a final object}

\msp{page 82}
\textit{and let us order $I$ by inclusion. Then $I$ is large relative to $\pi$, moreover the canonical functor $I \longrightarrow J$ $(J' \longmapsto e_{J'} =$ final object of $J')$ is cofinal} (\footnote{Introduce $I(J)$ before the statement of the proposition. \textbf{N.B.} The cofinality condition $I(J) \longrightarrow J$ is not automatic, as we see by taking $J = B_G$, $G = \{\pm 1\}$.},\footnote{\textbf{N.B.} The cofinality hypothesis is superfluous if $J$ is rigid, i.e. every endomorphism of an object of $J$ is the identity.},\footnote{\ Ed. note.\ \ In this condition (d), the ``\textit{moreover the functor \dots\ is cofinal}'' is added by Grothendieck \textit{a posteriori}, which explains his vain attempt to establish it in the implication d $\Longrightarrow$ a below, an implication which becomes obvious after this addition. On the other hand, conditions (c) and (d) are \textit{not} implied by the equivalent conditions (a) and (b). For example, if $J$ is the category having a single object and a single nonidentity morphism $p$ such that $p^2 = p$, then $J$ satisfies conditions (a) and (b), but \textit{not} conditions (c) and (d).}).
\end{itemize}
(\footnote{Other conditions:
\begin{itemize}[leftmargin=3em]
\item[e)] For every set $A$ with $\card A \leqslant \pi$, $\mathrm{diag}_J : J \longrightarrow J^A$ is cofinal.
\item[e$'$)] (If $A_0$ is a given set, of cardinality $= \pi$.) $J \longrightarrow J^{A_0}$ is cofinal.
\end{itemize}
See also 4.9.6.9.1 and 4.9.6.9.2, pages 237, 238.})

\textsc{Proof}.
\begin{itemize}[leftmargin=7em,labelwidth=6em,align=left]
\item[a $\;\Longrightarrow\;$ b] immediate verification, using the criterion F 1, F 2 for cofinality for $I$ filtering.
\item[b $\;\Longrightarrow\;$ c] by AQT \texttt{[âne qui trotte: at a donkey's trot, i.e. by routine verification]}.
\item[c $\;\Longrightarrow\;$ d] by AQT \texttt{[âne qui trotte: at a donkey's trot, i.e. by routine verification]}.
\end{itemize}
% foot: Version 17 décembre 2025

\printedpage{18}
% head: A. Grothendieck : Dérivateurs, Ch. XVIII, Édition des Universités de Montpellier II et Paris VII
\begin{itemize}[leftmargin=7em,labelwidth=6em,align=left]
\item[d $\;\Longrightarrow\;$ a] Let us consider the ordered set large relative to $\pi$ constructed in the statement of d), and the \textit{functor}
\[
I \;\longrightarrow\; J
\]
given by $J' \longmapsto e_{J'}$ $=$ final object of $J'$. (Verification that it is a functor immediate.) This functor is surjective, hence satisfies the cofinality condition F 1. Condition F 2 also follows from d) by looking at the three objects $J'$, $j \xrightarrow{\;u_1\;} e_{J'}$, $j \xrightarrow{\;u_2\;} e_{J'}$ of $I$,
\[
\begin{tikzcd}[column sep=small, row sep=small]
J', & e_{J'} \\
j \arrow[ur, shift left, "u_1"] \arrow[ur, shift right, "u_2"'] &
\end{tikzcd}
\]
and by taking a $J''$ which is an upper bound for all three, then $e_{J'} \longrightarrow e_{J''}$ equalizes $u_1$, $u_2$ (\footnote{A bit brief, the case where $e_{J'} = j$ escapes this argument, and requires that one \texttt{[incomplete sentence]}}).
\end{itemize}

% margin: vertical bar in the left margin beside Corollaire 4.4.5
\textbf{Corollary 4.4.5.} \textit{Let $F \in \Ob C^{\wedge}$. For $F$ to belong to} $\Ind'_\pi(C)$\textit{, it is necessary and sufficient that $C/F$ satisfy the set of four equivalent conditions of 4.4.4.}

\textbf{Corollary 4.4.6.} \textit{Let $J$ be a category stable under} $\varinjlim$ \textit{of cardinality $\leqslant \pi$. Then $J$ is large relative to $\pi$.}

Indeed, conditions $\mathrm{PF}_{\pi}\,1$, $\mathrm{PF}_{\pi}\,2$ are visibly satisfied.

\msp{page 83}
\textbf{4.4.6.} I give myself a functor % [sic] 4.4.6. -> 4.4.7. (?), numéro déjà employé par le Corollaire 4.4.6 qui précède
\[
J \xrightarrow{\;\varphi\;} \Ind'_\pi C \;, \tag*{(4.4.6.1)}
\]
$J$ a small category, and we look for a $\varinjlim$ in $\Ind'_\pi C$, postulating if need be the existence of $\varinjlim$ of cardinality $\leqslant \pi$ in $C$, and the hypothesis $\card\Fl J \leqslant \pi$. The most natural procedure for reducing the question of the existence of $\varinjlim$ in $\Ind'_\pi C$, to the similar question in $C$, is made explicit in SGA 4 I 8.8.3. We consider the fully faithful inclusion
\[
C \;\hookxrightarrow{\;\alpha\;}\; \Ind'_\pi C \;, \tag*{(4.4.6.2)}
\]
which commutes with $\varinjlim$ of cardinality $\leqslant \pi$ (those which are representable in $C$, cf. p. 78 and reference to SGA 4 I 9.8), whence
\[
\uHom(J,C) \;\hookxrightarrow{\;\alpha\;}\; \uHom(J,\Ind'_\pi C) \tag*{(4.4.6.3)}
\]
also fully faithful, and commuting with $\varinjlim$ of cardinality $\leqslant \pi$ (representable in the first member). This arrow fits into the commutative square relative to the datum of a $j \in \Ob J$, and of the inclusion functors
\[
e \xrightarrow{\;\varepsilon_j\;} J \;, \qquad e_j(e) \;=\; j \;; \tag*{(4.4.6.4)} % [sic] e_j(e) -> \varepsilon_j(e), la flèche est notée \varepsilon_j
\]
% foot: Version 17 décembre 2025

\end{document}
